\documentclass[12pt,a4paper]{exam}
\usepackage[utf8]{inputenc}
\usepackage{amsmath,amssymb,amsfonts}
\usepackage{geometry}
\usepackage{enumitem}
\usepackage{graphicx}
\usepackage{hyperref}
\usepackage{xcolor}
\geometry{a4paper, top=2cm, bottom=2cm, left=2.5cm, right=2.5cm}
\hypersetup{colorlinks=true, linkcolor=blue, urlcolor=blue}
\renewcommand{\thequestion}{\textbf{\arabic{question}}}
\renewcommand{\questionlabel}{\thequestion.}
\pointname{ marks}
\pointsdroppedatright
\bracketedpoints
\setlength{\parindent}{0pt}
\setlength{\emergencystretch}{3em}
\allowdisplaybreaks
\newsavebox{\schmathbox}
\newcommand{\fitmath}[1]{%
  \sbox{\schmathbox}{$\displaystyle #1$}%
  \ifdim\wd\schmathbox>\linewidth
    \resizebox{\linewidth}{!}{\usebox{\schmathbox}}%
  \else
    \usebox{\schmathbox}%
  \fi}

\begin{document}

\thispagestyle{empty}
\begin{center}
  \textbf{\LARGE Diag} \\[0.3cm]
  \textbf{Time Allowed: 3 Hours} \hfill \textbf{Maximum Marks: 80} \\[0.3cm]
  \rule{\textwidth}{0.5pt}
\end{center}

\vspace{0.3cm}

\noindent\textbf{General Instructions:}
\begin{enumerate}
  \item {\normalfont\normalcolor This question paper contains 40 questions. All questions are compulsory.}
  \item {\normalfont\normalcolor Questions are grouped into sections A through E.}
  \item {\normalfont\normalcolor Use of calculators is NOT permitted.}
\end{enumerate}
\bigskip
\noindent\textbf{\large Section A: Multiple Choice \& Assertion-Reason (1 mark each)}

\begin{questions}
\question[1] Two concentric circles are of radii $5$ cm and $3$ cm. The length of the chord of the larger circle which touches the smaller circle is:
\begin{enumerate}[label=(\Alph*)]
  \item (A) 6 cm
  \item (B) 8 cm
  \item (C) 10 cm
  \item (D) 7 cm
\end{enumerate}
\vspace{0.15cm}

\question[1] If $TP$ and $TQ$ are two tangents to a circle with center $O$ so that $\angle POQ = 110^\circ$, then $\angle PTQ$ is equal to:
\begin{enumerate}[label=(\Alph*)]
  \item $(A) 60^\circ$
  \item $(B) 80^\circ$
  \item $(C) 70^\circ$
  \item $(D) 90^\circ$
\end{enumerate}
\vspace{0.15cm}

\question[1] A tower stands vertically on the ground. From a point on the ground, which is $15$ m away from the foot of the tower, the angle of elevation of the top of the tower is found to be $60^\circ$. The height of the tower is:
\begin{enumerate}[label=(\Alph*)]
  \item (A) $5\sqrt{3}$ m
  \item (B) $10\sqrt{3}$ m
  \item (C) $15\sqrt{3}$ m
  \item (D) $20\sqrt{3}$ m
\end{enumerate}
\vspace{0.15cm}

\question[1] A kite is flying at a height of $60$ m above the ground. The string attached to the kite is temporarily tied to a point on the ground. The inclination of the string with the ground is $30^\circ$. The length of the string is:
\begin{enumerate}[label=(\Alph*)]
  \item (A) $60$ m
  \item (B) $30$ m
  \item (C) $90$ m
  \item (D) $120$ m
\end{enumerate}
\vspace{0.15cm}

\question[1] The angle of elevation of the top of a pole from a point on the ground $20$ m away from the foot of the pole is $45^\circ$. The height of the pole is:
\begin{enumerate}[label=(\Alph*)]
  \item (A) $10$ m
  \item (B) $20$ m
  \item (C) $20\sqrt{3}$ m
  \item (D) $40$ m
\end{enumerate}
\vspace{0.15cm}

\question[1] If the length of the shadow of a tower is $\sqrt{3}$ times the height of the tower, then the angle of elevation of the sun is:
\begin{enumerate}[label=(\Alph*)]
  \item (A) $30^\circ$
  \item (B) $45^\circ$
  \item (C) $60^\circ$
  \item (D) $90^\circ$
\end{enumerate}
\vspace{0.15cm}

\question[1] An observer $1.5$ m tall is $28.5$ m away from a chimney. The angle of elevation of the top of the chimney from her eyes is $45^\circ$. The height of the chimney is:
\begin{enumerate}[label=(\Alph*)]
  \item (A) $28.5$ m
  \item (B) $27$ m
  \item (C) $30$ m
  \item (D) $31.5$ m
\end{enumerate}
\vspace{0.15cm}

\question[1] Assertion (A): The length of a tangent drawn from an external point to a circle is always equal to the radius of the circle. Reason (R): Tangents drawn from an external point to a circle are equal in length.
\begin{enumerate}[label=(\Alph*)]
  \item (A) Both Assertion (A) and Reason (R) are true and Reason (R) is the correct explanation of Assertion (A)
  \item (B) Both Assertion (A) and Reason (R) are true but Reason (R) is NOT the correct explanation of Assertion (A)
  \item (C) Assertion (A) is true but Reason (R) is false
  \item (D) Assertion (A) is false but Reason (R) is true
\end{enumerate}
\vspace{0.15cm}

\question[1] Assertion (A): If the angle between two tangents drawn from a point P to a circle of radius $r$ and center O is $60^\circ$, then $OP = 2r$. Reason (R): The line segment joining the center to the external point bisects the angle between the two tangents.
\begin{enumerate}[label=(\Alph*)]
  \item (A) Both Assertion (A) and Reason (R) are true and Reason (R) is the correct explanation of Assertion (A)
  \item (B) Both Assertion (A) and Reason (R) are true but Reason (R) is NOT the correct explanation of Assertion (A)
  \item (C) Assertion (A) is true but Reason (R) is false
  \item (D) Assertion (A) is false but Reason (R) is true
\end{enumerate}
\vspace{0.15cm}

\question[1] Assertion (A): A tangent to a circle is perpendicular to the radius at the point of contact. Reason (R): The perpendicular distance from the center of the circle to the tangent is equal to the radius of the circle.
\begin{enumerate}[label=(\Alph*)]
  \item (A) Both Assertion (A) and Reason (R) are true and Reason (R) is the correct explanation of Assertion (A)
  \item (B) Both Assertion (A) and Reason (R) are true but Reason (R) is NOT the correct explanation of Assertion (A)
  \item (C) Assertion (A) is true but Reason (R) is false
  \item (D) Assertion (A) is false but Reason (R) is true
\end{enumerate}
\vspace{0.15cm}

\question[1] Assertion (A): The number of tangents that can be drawn from a point inside a circle is zero. Reason (R): A tangent to a circle intersects the circle at only one point.
\begin{enumerate}[label=(\Alph*)]
  \item (A) Both Assertion (A) and Reason (R) are true and Reason (R) is the correct explanation of Assertion (A)
  \item (B) Both Assertion (A) and Reason (R) are true but Reason (R) is NOT the correct explanation of Assertion (A)
  \item (C) Assertion (A) is true but Reason (R) is false
  \item (D) Assertion (A) is false but Reason (R) is true
\end{enumerate}
\vspace{0.15cm}

\question[1] Assertion (A): If the distance between two parallel tangents of a circle of radius $3$ cm is $6$ cm. Reason (R): Parallel tangents touch the circle at the ends of the diameter.
\begin{enumerate}[label=(\Alph*)]
  \item (A) Both Assertion (A) and Reason (R) are true and Reason (R) is the correct explanation of Assertion (A)
  \item (B) Both Assertion (A) and Reason (R) are true but Reason (R) is NOT the correct explanation of Assertion (A)
  \item (C) Assertion (A) is true but Reason (R) is false
  \item (D) Assertion (A) is false but Reason (R) is true
\end{enumerate}
\vspace{0.15cm}

\question[1] Assertion (A): The angle of elevation of the top of a tower from a point on the ground, which is $30 \text{ m}$ away from the foot of the tower of height $30 \sqrt{3} \text{ m}$, is $60^\circ$. 
Reason (R): In a right-angled triangle, $\tan(\theta) = \frac{\text{Opposite}}{\text{Adjacent}}$.
\begin{enumerate}[label=(\Alph*)]
  \item (A) Both Assertion (A) and Reason (R) are true and Reason (R) is the correct explanation of Assertion (A)
  \item (B) Both Assertion (A) and Reason (R) are true but Reason (R) is NOT the correct explanation of Assertion (A)
  \item (C) Assertion (A) is true but Reason (R) is false
  \item (D) Assertion (A) is false but Reason (R) is true
\end{enumerate}
\vspace{0.15cm}

\question[1] Assertion (A): If the length of the shadow of a vertical pole is equal to its height, then the angle of elevation of the Sun is $45^\circ$. 
Reason (R): The angle of elevation is defined only when the object is taller than the observer.
\begin{enumerate}[label=(\Alph*)]
  \item (A) Both Assertion (A) and Reason (R) are true and Reason (R) is the correct explanation of Assertion (A)
  \item (B) Both Assertion (A) and Reason (R) are true but Reason (R) is NOT the correct explanation of Assertion (A)
  \item (C) Assertion (A) is true but Reason (R) is false
  \item (D) Assertion (A) is false but Reason (R) is true
\end{enumerate}
\vspace{0.15cm}

\question[1] Assertion (A): The angle of depression of an object from a point $h$ meters above the ground is always equal to the angle of elevation of that point from the object on the ground. 
Reason (R): The angle of elevation and depression are alternate interior angles for parallel horizontal lines.
\begin{enumerate}[label=(\Alph*)]
  \item (A) Both Assertion (A) and Reason (R) are true and Reason (R) is the correct explanation of Assertion (A)
  \item (B) Both Assertion (A) and Reason (R) are true but Reason (R) is NOT the correct explanation of Assertion (A)
  \item (C) Assertion (A) is true but Reason (R) is false
  \item (D) Assertion (A) is false but Reason (R) is true
\end{enumerate}
\vspace{0.15cm}

\question[1] Assertion (A): If a ladder of length $10 \text{ m}$ reaches a wall at a height of $5 \text{ m}$, the angle made by the ladder with the ground is $30^\circ$. 
Reason (R): $\sin(\theta) = \frac{\text{Opposite}}{\text{Hypotenuse}}$.
\begin{enumerate}[label=(\Alph*)]
  \item (A) Both Assertion (A) and Reason (R) are true and Reason (R) is the correct explanation of Assertion (A)
  \item (B) Both Assertion (A) and Reason (R) are true but Reason (R) is NOT the correct explanation of Assertion (A)
  \item (C) Assertion (A) is true but Reason (R) is false
  \item (D) Assertion (A) is false but Reason (R) is true
\end{enumerate}
\vspace{0.15cm}

\question[1] Assertion (A): As an observer moves closer to a tower, the angle of elevation of its top increases. 
Reason (R): $\tan(\theta)$ increases as $\theta$ increases from $0^\circ$ to $90^\circ$.
\begin{enumerate}[label=(\Alph*)]
  \item (A) Both Assertion (A) and Reason (R) are true and Reason (R) is the correct explanation of Assertion (A)
  \item (B) Both Assertion (A) and Reason (R) are true but Reason (R) is NOT the correct explanation of Assertion (A)
  \item (C) Assertion (A) is true but Reason (R) is false
  \item (D) Assertion (A) is false but Reason (R) is true
\end{enumerate}
\vspace{0.15cm}

\end{questions}
\bigskip
\noindent\textbf{\large Section B: Very Short Answer (2 marks each)}

\begin{questions}
\question[2] If a tangent $PA$ and $PB$ from a point $P$ to a circle with center $O$ are inclined to each other at an angle of $80^{\circ}$, then find $\angle POA$.
\vspace{0.15cm}

\question[2] A quadrilateral $\text{ABCD}$ is drawn to circumscribe a circle. If $AB = 6\text{ cm}$, $BC = 7\text{ cm}$, and $CD = 4\text{ cm}$, find the length of $AD$.
\vspace{0.15cm}

\question[2] From a point $Q$, the length of the tangent to a circle is $24\text{ cm}$ and the distance of $Q$ from the center is $25\text{ cm}$. Find the radius of the circle.
\vspace{0.15cm}

\question[2] Two concentric circles are of radii $5\text{ cm}$ and $3\text{ cm}$. Find the length of the chord of the larger circle which touches the smaller circle.
\vspace{0.15cm}

\question[2] If the angle between two tangents drawn from an external point $P$ to a circle of radius $r$ and center $O$ is $60^{\circ}$, find the length of $OP$.
\vspace{0.15cm}

\question[2] A ladder $15$ m long makes an angle of $60^\circ$ with the wall. Find the height of the point where the ladder touches the wall.
\vspace{0.15cm}

\question[2] The angle of elevation of the top of a tower from a point on the ground, which is $30$ m away from the foot of the tower, is $30^\circ$. Find the height of the tower.
\vspace{0.15cm}

\question[2] A kite is flying at a height of $60$ m above the ground. The string attached to the kite is temporarily tied to a point on the ground. The inclination of the string with the ground is $60^\circ$. Find the length of the string.
\vspace{0.15cm}

\question[2] A $1.5$ m tall boy is standing at some distance from a $30$ m tall building. The angle of elevation from his eyes to the top of the building increases from $30^\circ$ to $60^\circ$ as he walks towards the building. Find the distance he walked towards the building.
\vspace{0.15cm}

\question[2] From a point on the ground, the angles of elevation of the bottom and top of a transmission tower fixed at the top of a $20$ m high building are $45^\circ$ and $60^\circ$ respectively. Find the height of the tower.
\vspace{0.15cm}

\end{questions}
\bigskip
\noindent\textbf{\large Section C: Short Answer (3 marks each)}

\begin{questions}
\question[3] Two concentric circles are of radii $5 \text{ cm}$ and $3 \text{ cm}$. Find the length of the chord of the larger circle which touches the smaller circle.
\vspace{0.15cm}

\question[3] A quadrilateral ABCD is drawn to circumscribe a circle. Prove that $AB + CD = AD + BC$.
\vspace{0.15cm}

\question[3] If a tangent $PA$ and $PB$ from a point $P$ to a circle with centre $O$ are inclined to each other at an angle of $80^{\circ}$, find $\angle POA$.
\vspace{0.15cm}

\question[3] Prove that the parallelogram circumscribing a circle is a rhombus.
\vspace{0.15cm}

\question[3] A point $P$ is $13 \text{ cm}$ from the centre of a circle. The length of the tangent drawn from $P$ to the circle is $12 \text{ cm}$. Find the radius of the circle.
\vspace{0.15cm}

\question[3] A tower stands vertically on the ground. From a point on the ground, which is $15$ m away from the foot of the tower, the angle of elevation of the top of the tower is found to be $60^\circ$. Find the height of the tower.
\vspace{0.15cm}

\question[3] A $1.5$ m tall boy is standing at some distance from a $30$ m tall building. The angle of elevation from his eyes to the top of the building increases from $30^\circ$ to $60^\circ$ as he walks towards the building. Find the distance he walked towards the building.
\vspace{0.15cm}

\question[3] From the top of a $7$ m high building, the angle of elevation of the top of a cable tower is $60^\circ$ and the angle of depression of its foot is $45^\circ$. Determine the height of the tower.
\vspace{0.15cm}

\question[3] A kite is flying at a height of $60$ m above the ground. The string attached to the kite is temporarily tied to a point on the ground. The inclination of the string with the ground is $60^\circ$. Find the length of the string, assuming that there is no slack in the string.
\vspace{0.15cm}

\end{questions}
\bigskip
\noindent\textbf{\large Section D: Long Answer (4-5 marks each)}

\begin{questions}
\question[5] Prove that the lengths of tangents drawn from an external point to a circle are equal.
\vspace{0.15cm}

\question[5] Two concentric circles are of radii $5\text{ cm}$ and $3\text{ cm}$. Find the length of the chord of the larger circle which touches the smaller circle.
\vspace{0.15cm}

\question[5] A quadrilateral $\text{ABCD}$ is drawn to circumscribe a circle. Prove that $AB + CD = AD + BC$.
\vspace{0.15cm}

\question[5] If a hexagon $\text{ABCDEF}$ circumscribes a circle, prove that $AB + CD + EF = BC + DE + FA$.
\vspace{0.15cm}

\end{questions}
\bigskip

\end{document}
